\documentclass[11pt,a4paper]{article}

% Informações do Preambulo
\input{Pre_textual/preambulo.tex}


% Iniciando o documento
\begin{document}


% Cabeçalho do evento
\input{Pre_textual/cabecalho.tex}

\vspace{1cm}

% Título do trabalho
{\fontsize{12}{12}\selectfont
\begin{center}
% Inserir o título do trabalho abaixo
\textbf {TÍTULO (Times New Roman, 12, Negrito, Centralizado)}
\end{center}}

% Nome do(s) autor(es)
\begin{center}
  NOME COMPLETO DO AUTOR$^1$, AUTOR$^2$, AUTOR$^3$, AUTOR$^4$\\
 (Times New Roman, 11, Centralizado, Máximo quatro autores)
\end{center}

% Informações sobre o(s) autor(es)
\begingroup
\sloppy
{\fontsize{10}{10}\selectfont
\noindent{$^1$} Titulação, filiação, endereço de e-mail (Times New Roman, 10, Justificado) % Para e-mail utilize \href{email@email.br}.
\\
{$^2$}  Caso seja de outra instituição ou câmpus, diferente do autor 1.\\
{$^3$} Caso seja de outra instituição ou câmpus, diferente do autor 1.\\
{$^4$} Caso seja de outra instituição ou câmpus, diferente do autor 1.
}
\endgroup

% Deletar esse comentário antes de submeter o artigo
{\fontsize{10}{10}\selectfont
\noindent {\color{red}{\textit{** O trabalho será avaliado na modalidade “às cegas”, ou seja, na submissão do artigo, deve-se apagar os nomes dos autores e não deve ser feita nenhuma citação no texto do nome da instituição, câmpus, empresa ou pessoas vinculadas ao trabalho.}}
}}

% Definir a área de conhecimento
\begin{center}
{\fontsize{10}{10}\selectfont 
\textbf{Área de conhecimento} (Tabela CNPq): Área Exemplo – 0.00.00.00-0}
\end{center}

\vspace{0.25cm}

\noindent
{\fontsize{10}{10}\selectfont
\textbf{RESUMO:} O propósito destas instruções é orientar o(s) autor(es) quanto à formatação dos resumos expandidos e artigos completos a serem submetidos ao VI Encontro de Iniciação Científica e Tecnológica (VII EnICT) do Instituto Federal de Educação, Ciência e Tecnologia de São Paulo - Câmpus Araraquara. Os documentos devem ser redigidos de acordo com as normas para elaboração. O arquivo de submissão deverá estar desbloqueado no formato \textit{portable document format} (pdf) compatível com o Adobe Acrobat Reader$^{\text{TM}}$. A palavra RESUMO e o texto que segue devem iniciar na mesma linha. O texto deve ser claro, sucinto e, obrigatoriamente, explicar o(s) objetivo(s) pretendido(s), procurando justificar sua importância, os principais procedimentos adotados, os resultados mais expressivos e as conclusões, contendo, no máximo, 200 palavras. O RESUMO não deverá conter fórmulas, citações e referências bibliográficas. Ele deverá ser escrito em fonte Times New Roman, tamanho 10 e o texto deve estar justificado. Todas as instruções supracitadas devem ser apagadas antes da submissão do trabalho. O resumo expandido ou artigo completo apresentado no evento será publicado nos Anais.

\vspace{0.25cm}

\noindent
\textbf{PALAVRAS-CHAVE:} máximo de seis, separadas por ponto e vírgula (;), escritas em letras minúsculas, organizadas em ordem alfabética crescente. É desejável que não se repita palavras utilizadas no título.}


\vspace{0.5cm}
\noindent {\fontsize{12}{12}\selectfont
\textbf{INTRODUÇÃO}}

O resumo expandido deve ter no máximo 04 páginas e o artigo completo deve possuir no mínimo 05 e no máximo 06 páginas. Ambos podem ser escritos em Português ou Inglês, no formato de página A4, coluna única, margens superior, de 3 cm, inferior, 2,5 cm, esquerda e direita, de 2 cm e não é necessário numerar as páginas. O espaçamento entre linhas deve ser simples. Os títulos de seções devem ser escritos em fonte Times New Roman, tamanho 12 e devem estar justificados. O texto principal deve ser escrito em fonte Times New Roman, tamanho 11 e também deve estar justificado. O texto deve ser escrito em 3$^{\text{a}}$ pessoa, ou seja, na forma impessoal. As citações e as referências bibliográficas devem seguir as normas ABNT NBR10520 e ABNT6023, respectivamente.

Após a divulgação dos resultados, os autores dos trabalhos aceitos deverão fazer uma nova submissão com a versão final do texto e incluir os nomes dos autores. 

A introdução tem como objetivo posicionar o leitor dentro do tema que está sendo desenvolvido. Deve conter a delimitação do universo da pesquisa, problema, objetivos, justificativa, hipóteses, metodologia de trabalho e sua relevância. Evitar o uso de palavras e ideias repetidas. Seja objetivo, claro e conciso em suas considerações, evitando que a introdução do trabalho fique longa e cansativa.



\vspace{0.5cm}
\noindent
{\fontsize{12}{12}\selectfont
\textbf{FUNDAMENTAÇÃO TEÓRICA}}

Apresenta as referências nas quais se baseia a pesquisa. Pesquisar é procurar informação que não se sabe e precisa saber. Para isso o trabalho científico precisa de fontes fidedignas. É necessário conhecer as fontes e os métodos de busca.  De acordo com \citeonline[p. 77]{Bocchi:04} ``no trabalho científico, não há espaço para o senso comum ou a falta de certeza: o pesquisador não acha nada, ele se apoia em teorias científicas, observações empíricas e elaborações estatísticas."

Todas as etapas na construção de um documento científico dependem de uma revisão bibliográfica, resultando em um texto que é a fundamentação teórica. Elabora-se um texto a partir de material já publicado, constituído, principalmente, de: \begin{itemize}[noitemsep]
    \item Livros;
    \item Artigos em revistas/periódicos;
    \item Teses de doutorado;
    \item Dissertações de mestrado;
    \item Monografias em geral;
    \item Outros trabalhos como, por exemplo, documentários.
\end{itemize}
O referencial teórico tem também outras funções \cite{Vergara:97}:
\begin{itemize}[noitemsep]
    \item Permite que o autor tenha maior clareza;
    \item Facilita a formulação de hipóteses e de suposições;
    \item Sinaliza para o método mais adequado à solução do problema;
    \item Permite identificar qual o procedimento mais pertinente para a coleta e o tratamento dos dados, bem como o conteúdo do procedimento escolhido;
    \item Uma revisão de literatura serve para verificar o que já foi pesquisado do assunto;
    \item O autor pode mostrar suas preocupações e preferências, apontar para o leitor as lacunas que percebe na bibliografia ou as discordâncias que com ela tem ou os pontos que considera que precisam ser confirmados.
\end{itemize}

\vspace{0.5cm}
\noindent
{\fontsize{12}{12}\selectfont
\textbf{METODOLOGIA}}

O desenvolvimento da metodologia deve ser pertinente ao tema, procurando alcançar os objetivos já apresentados. Dessa forma, espera-se uma descrição adequada dessa metodologia utilizada para desenvolver o trabalho. 

Caso use equações, numerar de acordo com o exemplo dado a seguir, descrevendo os parâmetros/variáveis. 

\noindent 
\textbf{Modelo de equação:}
\begin{equation}
T=\frac{-b\pm\sqrt{\Delta}}{2a}
\end{equation}
onde,\\
$a$ - coeficiente do termo que possui a incógnita ao quadrado $(x^2)$,\\
$b$ - coeficiente do termo que possui a incógnita simples $(x)$,\\
$\Delta$ - delta $(\Delta=b^2-4ac)$,\\
$c$ - coeficiente do termo independente.


\vspace{0.5cm}
\noindent
{\fontsize{12}{12}\selectfont
\textbf{RESULTADOS E DISCUSSÃO}}

Apresentar os resultados de maneira organizada e lógica, fornecendo ao leitor as informações mais representativas. É comum o uso de figuras e tabelas para ilustrar os dados apresentados de maneira a facilitar a compreensão dos resultados obtidos. É importante seguir uma sequência adequada, que induza o leitor a tirar suas conclusões sobre os resultados encontrados e manter uma constância na sequência nas descrições. Os resultados devem ser descritivos e objetivos \cite{DC:08}.

A discussão deve ser argumentativa e ter um caráter mais subjetivo, apresentando os pontos a favor ou contra as ideias do(s) autor(es) que estudaram o mesmo assunto. Pode-se também analisar, interpretar, criticar e comparar os resultados da sua pesquisa com os já existentes sobre o assunto na literatura citada, discutindo possíveis implicações, significados e razões para concordância ou discordância em relação a outros autores. A principal finalidade é mostrar as relações entre os fatos observados e demonstrar o significado dos resultados obtidos  \cite{DC:08}.

Figuras, tabelas e gráficos devem ser apresentados com tamanho, qualidade e detalhes suficientes para a interpretação e composição gráfica final. Citar a fonte abaixo da ilustração, tabela e gráfico e, caso seja do próprio autor, citar como ``Fonte: Elaborada pelo autor". As legendas e fontes das figuras, tabelas e gráficos devem ser escritos em fonte Times New Roman, tamanho 10 e o texto deve estar centralizado, em negrito e seguir o padrão de numeração sucessivo arábico. Centralizar as Figuras, Tabelas e Gráficos.

\textbf{Gráficos e figuras:} devem apresentar-se sem bordas; a legenda deve ser posicionada acima da figura ou gráfico. Usar imagens coloridas se for necessário destacar alguma informação da Figura ou Gráfico. 

\textbf{Tabelas:} : evitar tabelas extensas e dados supérfluos; adequar seus tamanhos ao espaço útil do papel e colocar, na medida do possível, apenas linhas contínuas horizontais. As legendas devem ser autoexplicativas, localizadas acima da tabela, como no exemplo a seguir.  

\noindent
\textbf{Modelo de Figura:}
\begin{table}[ht]
    \centering
    \label{fig1}
   \caption{Esquema computacional da visão estéreo.}
    \includegraphics[scale=1]{Figuras/figura_1.png}\\
    {\scriptsize  Fonte: \citeonline{RA:09}.}
\end{table}

\noindent
\textbf{Modelo de Tabela:}
\begin{table}[ht]
\centering
\begin{tabular}{|c|c|c|c|}
\hline
\cellcolor[HTML]{C0C0C0} &\cellcolor[HTML]{C0C0C0}\textbf{Tempo de Processamento} 
& \cellcolor[HTML]{C0C0C0}\textbf{Tempo de Comunicação}
&\cellcolor[HTML]{C0C0C0} \textbf{Tempo total} \\ \hline
\cellcolor[HTML]{C0C0C0} \textbf{GPU} & 10.4 ms  & 158 ms & 168 ms \\ \hline
\cellcolor[HTML]{C0C0C0}\textbf{CPU} & 2180 ms & - & 2180 ms \\ \hline
\end{tabular}\\[-3mm]
\begin{center}
  {\scriptsize  Fonte: \citeonline{WAR:09}.}    
\end{center}
\end{table}

\vspace{0.5cm}
\noindent
{\fontsize{12}{12}\selectfont
\textbf{CONCLUSÕES}}

Devem basear-se, exclusivamente, nos resultados do trabalho para artigos completos e, para resumos expandidos, nos resultados preliminares quando houver e nas discussões sobre as hipóteses iniciais do projeto. Ressaltar se os objetivos previstos inicialmente foram cumpridos. Aspectos metodológicos também devem ser abordados: dificuldades durante a pesquisa, resultados obtidos (artigos completos) e qual o significado dos mesmos dentro da área de pesquisa. Quando possível, inserir outros tipos de estudos que poderiam ser feitos utilizando o mesmo tema e/ou outros aspectos a serem abordados a partir dos resultados apresentados e, para resumos expandidos, as próximas etapas da pesquisa.

\vspace{1cm}
\noindent
{\fontsize{12}{12}\selectfont
\textbf{AGRADECIMENTOS (Opcional)}}

Inserir, após as conclusões, de maneira sucinta os agradecimentos. Se o projeto for financiado por alguma agência de fomento, citar a fonte.



\vspace{1cm}
\noindent
{\fontsize{12}{12}\selectfont
\textbf{REFERÊNCIAS}}

{\color{red}
As referências devem ser listadas em ordem alfabética. Veja alguns exemplos:}


\begingroup
\renewcommand{\section}[2]{}
\bibliography{bibliografia.bib}
\endgroup

\end{document}